\section{Grassmann map} \label{sec;Gmap}

In this section we let $n=1$ or $n=3$, $S^{2n} \subset \R^{2n+1}$ be the unit sphere with respect to the standard Euclidean metic, and let $\nabla$ be the Levi-Civita connection. Recall that at any point $x \in S^{2n}$ we have
\begin{gather*}
\CC^{2n+1} = \big( T^{0,1}_x \oplus (N_x \otimes \CC) \big) \oplus T^{1,0}_x.
\end{gather*}

\begin{defn}For any $J$ on $S^{2n}$, the induced map $P^\perp\colon S^{2n} \to \cG_{n+1,2n+1}$ is defined so that~$P^\perp(x)$ is the orthogonal projection of $\CC^{2n+1}$ onto $T^{0,1}_x \oplus (N_x \otimes \CC)$.
\end{defn}

Lebrun \cite{L} considers the closely related map which is given by orthogonal projection onto~$T^{0,1}_x$, and shows it is an embedding. We'll see below that by including the normal bundle in the new map above, the prohibited condition on~$S^6$ becomes ``strongly integrable'' rather than ``integrable and orthogonal''. This again is motivated by the normal component in equation~\eqref{eq;NablaSphere}.

Theorem \ref{thm;delbarPperp} and Corollary~\ref{cor;noStrongJ} below are very much inspired by the author's reading of~\cite{L}. First, we show this map is an immersion, which is sufficient for our purposes.

\begin{lem} For any $J$ on $S^{2n}$, the map $P^\perp\colon S^{2n} \to \cG_{n+1,2n+1}$ is an immersion.
\end{lem}

\begin{proof}If $X \in T_x S^{2n}$ then $d_X P \in T_{P(x)} \cG_{n+1,2n+1}$ is a non-zero matrix since, for the real normal unit vector $n_x$ to $S^{2n}$ at $x$, we have
\begin{gather*}
\big(d_XP^\perp\big)(n_x) = d_X\big(P^\perp(n_x)\big) - P^\perp d_X(n_x) = \big(I-P^\perp \big) d_X(n_x) = \big(I-P^\perp\big)(X).
\end{gather*}
This shows $d_X P$ is non-zero whenever $X$ is non-zero since $X \in T_xM$ is real, but $P^\perp(X) \in T^{0,1}$ is not real.
\end{proof}

Recall that the tangent algebra $m$ associated to $J$ and $\nabla$ is given by
\begin{gather*}
m(X,Y) = \nabla_{JX} (J)Y - J \nabla_X (J) Y,
\end{gather*}
and that $J $ is strongly integrable with respect to $\nabla$ if and only if $\nabla_X J$ is $J$-linear in the variable~$X$, i.e., $m$ vanishes.

For a map $f\colon (M,J) \to (N,K) $ between almost complex manifolds, we define
\begin{gather*}
\overline{\partial}_X f = d_{JX} f - K \circ d_X f.
\end{gather*}

\begin{thm} \label{thm;delbarPperp}
For any $J$ on $S^{2n}$, the map $P^\perp\colon S^{2n} \to \cG_{n+1,2n+1}$ induced by $J$ satisfies
\begin{gather*}
\overline{\partial}_X P^\perp (n_x) = 0, \\
\overline{\partial}_X P^\perp (Y+iJY) = \big( I - P^\perp\big) (J m(X,Y)),
\end{gather*}
where $m(X,Y) = \nabla_{JX} (J)Y - J \nabla_X (J) Y$ is the tangent algebra.

Furthermore, $\overline{\partial} P^\perp \equiv 0$ if and only if~$J $ is strongly integrable with respect to~$\nabla$, i.e., the tangent algebra $m$ vanishes.
\end{thm}

\begin{proof}Recall from Section \ref{sec;grass} that the complex structure on $T_P \cG_{n+1,2n+1}$ is given by multiplication by $+i$ on $\operatorname{Hom}\big(\operatorname{Im}\big(P^\perp\big), \operatorname{Ker}\big(P^\perp\big)\big)$, and by self-adjointness, is given by multiplication by $-i$ on $\operatorname{Hom}\big(\operatorname{Ker}\big(P^\perp\big),\operatorname{Im}\big(P^\perp\big)\big)$.

As in the previous lemma we have
\begin{gather*}
 \overline{\partial}_X P^\perp (n_x) = -i \big( d_{X+iJX} P^\perp \big) (n_x) = -i \big(I-P^\perp\big)(X+iJX) = 0,
\end{gather*}
and, for arbitrary $Z \in T^{0,1}_x$, we have
\begin{gather*}
\big( d_{X+iJX} P^\perp \big) (Z) = \big(I-P^\perp\big) d_{X+iJX} (Z)= \big(I-P^\perp\big) ( \nabla_{X+iJX} (Z) - \langle X+iJX, Z \rangle n_x ) \\
\hphantom{\big( d_{X+iJX} P^\perp \big) (Z)}{} = \big(I-P^\perp\big) ( \nabla_{X+iJX} (Z) ),
 \end{gather*}
 since $P^\perp(n_x) = n_x$. Now, if $Z= Y+iJY$, then the last expression is equal to
 \begin{gather*}
 \big(I -P^\perp\big) \big( \nabla_X Y - \nabla_{JX} (JY) + i ( \nabla_X(JY) + \nabla_{JX} Y ) \big) \\
 \qquad{} = \big(I -P^\perp\big) i \big( \nabla_X(JY) + \nabla_{JX} Y - J(\nabla_X Y - \nabla_{JX} (JY)) \big) \\
 \qquad {}= i \big(I -P^\perp\big) \big( J ( \nabla_{JX} (J) Y - J\nabla_X(J) Y ) \big) = i \big(I - P^\perp\big) ( J m(X,Y) ).
\end{gather*}
 So,
\begin{gather*}
\overline{\partial}_X P^\perp (Y+iJY) = -i \big( d_{X+iJX} P^\perp \big) (Y + iJY) = ( I - P^\perp) ( J m(X,Y) ),
\end{gather*}
as claimed. Since $ J m(X,Y)$ is real, the right hand side vanishes if and only if~$m$ vanishes, i.e., $J $~is strongly integrable with respect to $\nabla$.

The remainder of the proof, showing that $\overline{\partial} P^\perp \equiv 0$ when $m$ vanishes, is entirely formal since the restriction $dP^\perp\colon \big(T^{0,1}\big)^\perp \to T^{0,1}$ is the adjoint of the restriction $dP^\perp\colon T^{0,1} \to \big(T^{0,1}\big)^\perp$. Explicitly, for any $W + i JZ = \big(T^{0,1}\big)^\perp$ we have with respect to the complex inner product $\langle - , - \rangle$,
\begin{gather*}
\big\langle \big( d_{X-iJX} P^\perp \big) (W+iJZ) , Y+i JY \big\rangle
= \big\langle W+iJZ , \big( d_{X+iJX} P^\perp \big)( Y+i JY )\big\rangle = 0
\end{gather*}
since $P^\perp$ is always self adjoint. Also, for $U + iV \in \big(T^{0,1}\big)^\perp$, we have
\begin{gather*}
\big\langle \big( d_{X-iJX} P^\perp \big) (W+iJZ) , U+iJV \big\rangle = 0,
\end{gather*}
since $d_{X-iJX} P^\perp \colon \big(T^{0,1}\big)^\perp \to T^{0,1}$, because $P^\perp$ is a projection operator.
\end{proof}

Note that by Theorem~\ref{thm;Jintorth}, any integrable orthogonal $J$ on~$S^6$ is strongly integrable with respect to $\nabla$, whereas locally there are non-orthogonal almost complex structures on~$S^6$ which are strongly integrable with respect to a torsion free connection. For example, from Euclidean space we have those non-orthogonal almost complex structures which are covariantly constant. Nevertheless, we can conclude:

\begin{cor} \label{cor;noStrongJ}There are no almost complex structures on $S^6$ which are strongly integrable with respect to the Levi-Civita connection~$\nabla$ of the standard metric.
\end{cor}

\begin{proof} By the Theorem~\ref{thm;delbarPperp}, if $J$ is strongly integrable with respect to $\nabla$, then $P^\perp\colon S^6 \to \cG_{4,7}$ is a $J$-holomorphic immersion. Since $\cG_{4,7}$ is a K\"ahler manifold, this shows the pullback $\big(P^\perp \big)^*\omega$ is closed and non-degenerate, which contradicts $H^2\big(S^6\big) =0$.
\end{proof}
